\section{Discussion}\label{sec:discussion}

\subsection{What non-retention means, module by module}\label{sec:meaning}

A negative verdict from a scored test carries the weight of the test's power, and Section~\ref{sec:operating} lets that weight be assigned per module. For M1 the verdict is strong. At the parameters of the series on which it was tested, the rule retains a true autonomous Schaefer module in 97--99\% of replicates, so its non-retention on both specifications, with a persistence margin that excludes zero on Specification B, is evidence that an autonomous production model driven by mean historical catch does not describe this stock's year-to-year SSB transitions. Section~\ref{sec:periods} says why in plain terms: after 1992 the training-mean catch exceeds anything the stock could produce, and the module forecasts extinction at every origin for two decades. For M2, M3 and M4 the verdict is weak. Against a true stock-flow process at recovery-window parameters the rule has 9--11\% power, so non-retention is compatible with a catch-supplied production structure being present and too faint, at these noise levels and sample sizes, to clear a 5\% band at both horizons. The one-step skill of these modules in the 1996--2012 band of Specification A, and the shrinking of their one-step deficits under the consistent catch-timing convention, are consistent with that reading; their five-year deficits and their failure on Specification B are not evidence against it so much as evidence that what they carry is a short-memory growth term that does not extrapolate. For M1b the verdict is not about depensation at all: the parameter the module adds is unidentified in the fixed windows and takes implausible values in the rolling ones, so the test could not have retained a true depensatory module (power 0.5--1.5\%) and its non-retention says nothing about whether an Allee effect exists.

The processes outside the class qualify all of this. Under a drifting-productivity truth or a smooth path observed with error, the rule retains some module, usually M1, in most replicates. Had a module been retained on the cod series, the retention would have shown that something beyond persistence is present, not that the retained structure is that something. The rule is a filter against the absence of structure, with a declared and measured false-retention rate inside the class, and it should be reported as such.

\subsection{Why persistence is hard to beat on an assessed state}\label{sec:whypersist}

Three features of assessment output make persistence a demanding benchmark on SSB. The series is a retrospective state estimate, smoothed by the assessment model's process-error structure, so successive values are strongly dependent and the last value is an efficient summary of the recent past. The forecaster is fitted to the same smoothed series, and a production model fitted to slow transitions learns slow dynamics; it then adds a parametric shape and a catch plug that the smoothed series does not demand. And the quantity that moves the series most, the 1991--1995 collapse, is one that no member of the class can reproduce at the recorded catch, so the model carries the full persistence error there and adds its own. These are properties of scoring on the state, and they are the reason the paper scores on the state: hindcast cross-validation on the fitted index \citep{Kell2016,Kell2021,Carvalho2021} answers whether the assessment predicts its own data, which is a different and easier question than whether its dynamics predict next year's biomass. The finding that simple models are hard to beat on fish population series is not new \citep{Ward2014,Perretti2013}; what the present design adds is a rule that converts the comparison into a per-module decision with known operating characteristics, and a decomposition that shows where the comparison is won and lost.

\subsection{The collapse, the class and the reference point}\label{sec:classrp}

Proposition~\ref{prop:repeller} is elementary, but its consequences for how surplus-production forecasts are read are not always drawn. The lower equilibrium of a catch-driven Schaefer map is a product of the removals; it vanishes as the catch goes to zero and it exists in the compensatory branch without any Allee term. A forecaster of this class that shows a stock ``below the threshold from which it cannot recover'' is making a statement about the catch it has been given, and the same statement reverses when the catch is set to the post-moratorium level, which is why M2--M4 forecast a mechanical rebound from 1993 and why the autonomous modules do not. The collapse transitions themselves, with apparent production of $-60$ to $-240\kt$, are outside the class at any parameter value; the assessment reconciles them through unreported removals and a natural-mortality shift \citep{Cadigan2016}, and the debate over their causes \citep{Hutchings1994,Walters1996,Myers1997,Shelton1999} is a debate about which term outside the production equation absorbed them. The practical lesson is the one \citet{Walters1996} drew for assessment: a model that fits a plateau will forecast the plateau, and a scored forecast test run on the plateau years alone would have retained M1 with high confidence. The same map and the same catch-generated repeller underlie the viability analysis of the limit reference point in Abaee (2026, Robust viability of the 2J3KL limit reference point), where the question is not whether the map forecasts but what catch schedules keep the state above a threshold under worst-case disturbance; the obstruction that Proposition~\ref{prop:repeller} states for one dimension is treated in general form, for viability under incomplete observation, in (Abaee, 2026, An obstruction calculus). Nothing in the present test depends on those results.

\subsection{A regime signal and what it would take to register it}\label{sec:regimesignal}

The period decomposition is post hoc and is presented as description. It nonetheless points to a testable claim: that on a low-biomass, low-catch regime a catch-supplied production model has one-step skill over persistence of the order of 30--40\%, carried by a linear growth term with $r\approx 0.25$--$0.5$. A prospective test of that claim would fix the regime boundary in advance from a criterion external to the scores (for instance, the moratorium year plus the minimum training length), use the consistent catch-timing convention, apply the same band and both horizons, and require the result on both specifications. Specification B already supplies the warning: the 2013--2023 rebuild, which the Specification A series does not contain, is forecast by the same modules as a rise several hundred kilotonnes too large at five years, because a linear growth term fitted to 1995--2012 does not know where the rebuild will slow. A regime-conditional forecaster earns its retention only if it also knows when the regime ends, and the present results give no evidence that any member of the ladder does.

\subsection{The verdict layer}\label{sec:verdictlayer}

The RMSE comparison never uses a threshold, but every normative reading of a forecast does, and Section~\ref{sec:verdict} shows that the threshold layer moved more than the forecast layer after 2015. The same xteNCAM file is below the limit in 26 years under $\Blim$ and in 55 consecutive years under the 1980s-mean LRP; the LRP itself differs by 94.6\,kt between two vintages of the same model; and the 2023 revision changed the stock's reported zone without a change in the stock \citep{DFO2024a,Regular2025}. A forecast test can be run on either series, and this paper runs it on both; what it cannot do is score a forecast against a status that is a function of which series and which threshold the authority adopts. The 2022--2023 survey suspension adds the boundary case: for two assessment cycles the series the test presumes did not exist, the removals rule was rolled over unchanged, and the vessel that carries the series was condemned (SI S12). A review cadence that outruns the loss of the information it consumes is the condition analysed in Abaee (2026, Periodic review as sampled governance); the episode is its empirical instance, and it marks the practical limit of any scoring exercise that assumes a complete annual state.

The rule itself is portable. Applied under identical hurdles to a one-pool water-balance forecaster of Edwards Aquifer head in Abaee (2026, Does a one-pool water-balance model improve forecasts of Edwards Aquifer head), it withheld retention from a module that beat persistence by 2.96\% at one step and lost at five (SI S10). There the band and the two-horizon requirement decided; here the ranking alone decides. The two cases share the rule and nothing else, and no verdict is transferred between them.

\subsection{Limitations}\label{sec:limits}

The test is on one stock, two series, one model class and two horizons. Each specification uses a single retrospective vintage, which favours the models; the 1983--2015 overlap between vintages differs by 24.7\,kt on average and by 131\,kt in one year, so a real-time evaluation on contemporaneous assessments would be harsher on every forecaster, persistence included. Catch enters as a point series although both assessments treat it as bounded and the collapse years are known to be under-reported. The forecasts are point forecasts and the scores are RMSE and MAE; a probabilistic evaluation with proper scoring rules on predictive distributions \citep{HyndmanAthanasopoulos2021} would be the natural extension and would let the published xteNCAM intervals enter the comparison. Sample sizes are small on Specification A, the bootstrap intervals are wide, and the period bands were chosen after seeing the data. The archived runner's catch-timing offset affects the catch-supplied modules' one-step margins by 20--75\,kt without affecting the verdict; the archived numbers are the ones reported and the corrected ones are in SI S14. The design is a fixed computational protocol rather than a prospective registration: the tie band and the comparator declarations were formalised after the first scores were seen, and the variant passes were added in sequence (SI S1). The M1b optimum on Specification B is sensitive to the numerical environment at the $\pm17\kt$ level, immaterial to the verdict but a further sign that the depensatory branch is unidentified (SI S13). Finally, the operating characteristics were measured against truths inside the class and two specific truths outside it; power against other misspecified truths is unknown, and the figures reported are upper bounds for the in-class alternatives.
